\documentclass[a4paper]{article}
\usepackage[pdftex]{graphicx}
\usepackage[utf8]{inputenc}
\usepackage{enumerate}
\usepackage{href-ul}
\hypersetup{
	colorlinks=true,
	linkcolor=blue,
	urlcolor=blue}
\usepackage{siunitx}
\sisetup{locale=DE}
\usepackage{icomma}
\usepackage{amssymb}
\usepackage{geometry}
\geometry{a4paper, top=15mm, left=15mm, right=15mm, bottom=15mm,
	headsep=10mm, footskip=12mm}

\begin{document}
	
	\begin{enumerate}[1.]
		\item {\bf Mass and Inertia}
		
		\begin{enumerate}[a)]
			\item Under what conditions does a body behave sluggishly? What is meant by sluggish behavior?
			\item Although Thomas pedals with all his strength on a flat stretch of road, he travels at a constant speed. Explain why this is.
			\item Snow on your shoes can be removed by stomping the ground vigorously. Explain this matter from a physics perspective.
		\end{enumerate}
		
		
		\item {\bf Lever; Machinery} \\
		A symmetrically mounted lever is moved with a cable winch over a cable and roller combination. The attached weight G is $F_G=\SI{1.50}{\kilo\newton}$. At the lever the following applies: $a_1=\SI{1.50}{\meter}$ and $a_2=\SI{2.50}{\meter}$.
		\begin{enumerate}[a)]
			\item With what force does the winch have to pull the rope to ensure equilibrium?
			\item The cable winch consists of a crank with the radius $ R=\SI{40.0}{\centi\meter}$ and the cable drum with the radius $ r=\SI{10.0}{\centi\meter} $. \\
			Calculate what force $F_K$ must act on the crank in this case
			\item The rope anchorage at point A can withstand a maximum of 600 N. How far can the point of application of the weight G on the left lever arm at most be away from the lever axis under this condition?
		\end{enumerate}
		
		\includegraphics[width=12 cm]{ Hebel119.pdf}
		
		\item {\bf Torque}
		
		\begin{enumerate}[a)]
			\item The handle of a pair of pliers is moved $a_1=\SI{15}{\centi\meter}$ away from the pivot point with a force of
			$F_1=\SI{180}{\newton}$ compressed to cut a nail $a_2=\SI{3.5}{\centi\meter}$ from the pivot point.\\
			What force acts on the nail?
			\item Perpendicular to a symmetrically mounted lever rod, a force of $F_1=\SI{6.0}{\newton}$ acts on the left side $a_1=\SI{3.0}{\centi\meter}$ away from the pivot point. And on the right side $a_2=\SI{8,0}{\centi\meter}$ away from the pivot point a force of $F_2=\SI{3,0}{\newton}$ \\
			Make a sketch and add a third force so that balance prevails. (Reason!)
		\end{enumerate}
		
		\item {\bf Pulley}
		
		\begin{minipage}{0.45\textwidth}
			\begin{enumerate}[a)]
				\item How large does the force have to be at point C for equilibrium to prevail?
				\item How large are the forces at the suspension points A and B?
			\end{enumerate}
			\vfill
		\end{minipage}
		\hfill
		\begin{minipage}{0.5\textwidth}
			\includegraphics[width=5 cm]{flaschzug121.pdf}
		\end{minipage}
		
	\end{enumerate}
	
	\fbox{
		\begin{minipage}{0.5\textwidth}
			For the solution please \href{https://www.okuyakl.de/phys/p8mecL406/le406.pdf}{click here} or scan the QR code.\\
			Additional worksheets are available at
			
			\href{http://www.okuyakl.com}{www.okuyakl.com}
		\end{minipage}
		\hfill
		\begin{minipage}{0.4\textwidth}
			\includegraphics[width=1.5 cm]{../../viecher/zwe03}
			\includegraphics[width=3 cm]{qre406}
			\includegraphics[width=2 cm]{../../viecher/afanticon1}
			
	\end{minipage}}
	
\end{document}